% Generated draft: requires mathematical and source review.
\documentclass[11pt]{article}
\usepackage[margin=25mm]{geometry}
\usepackage{fontspec}
\setmainfont{Latin Modern Roman}
\usepackage{amsmath,amssymb,mathtools,unicode-math}
\setmathfont{Latin Modern Math}
\usepackage{enumitem,needspace,xurl,hyperref}
\setlength{\parindent}{0pt}
\setlength{\parskip}{5pt}
\setlength{\emergencystretch}{3em}
\newcommand{\N}{\mathbb{N}}
\newcommand{\Z}{\mathbb{Z}}
\newcommand{\Q}{\mathbb{Q}}
\newcommand{\R}{\mathbb{R}}
\newcommand{\C}{\mathbb{C}}
\newcommand{\F}{\mathbb{F}}
\newcommand{\A}{\mathbb{A}}
\newcommand{\PP}{\mathbb{P}}
\newcommand{\E}{\mathbb{E}}
\newcommand{\eps}{\varepsilon}
\newcommand{\dd}{\,\mathrm{d}}
\begin{document}
\section*{UnsolvedMath ID: 2307016}
\textbf{Source problem code:} catalogue entry
\subsection{Definitions and mathematical statement}
Let $\gamma$ be a Jordan arc in $\C$ and let $\mu$ be a measure supported on $\gamma$. The source considers the Laplace transform
\[
f(z)=\int_\gamma e^{z\zeta}\,d\mu(\zeta).
\]
In the usual compact-arc, finite-Borel-measure setting, this defines an entire function; for a complex measure the sufficient bound is
\[
|f(z)|\leq |\mu|(\gamma)\exp\!\left(|z|\max_{\zeta\in\gamma}|\zeta|\right),
\]
where $|\mu|$ is total variation. These finiteness and measure-type conventions are not stated in the supplied record, so they are explanatory conditions rather than confirmed source assumptions.

J. Korevaar asks whether every transform of the intended class has ``asymptotically analytic'' growth as $|z|\to\infty$. The source explicitly allows that the answer might be no. It then asks whether, even if that growth assertion fails, the zeros of $f$ might always have ``measurable distribution'' in A. Pfluger's sense.

Neither quoted technical property is defined in the supplied record. They cannot safely be replaced by complete regular growth, a particular indicator-function asymptotic, or a chosen angular zero-density limit without the original definitions. For a nonzero entire $f$, zeros can be encoded by the counting measure
\[
\nu_f=\sum_{f(a)=0}\operatorname{ord}_a(f)\,\delta_a,
\]
but the normalization, limiting process, and angular conditions meant by Pfluger's measurable distribution remain unspecified. The identically zero transform must also be excluded or treated separately before discussing discrete zero distribution. The supplied summary reports openness as of the 2018 edition and no definitive resolution located.
\subsection{Short English statement}
Korevaar asks first for a precise growth property of Laplace transforms of measures on Jordan arcs, and then for a potentially weaker zero-distribution property. The technical meanings, measure hypotheses, and treatment of the zero transform are missing, so neither property is replaced by a guessed standard notion.
\subsection{Sources}
\begin{itemize}
\item Local UnsolvedMath registry, permanent ID 2307016, exact source code catalogue entry, ``Research Problems in Function Theory --- Problem 7.16''; status open, classification OPEN-TRIAGE.
\item Hayman, \emph{Research Problems in Function Theory} (2018), Problem 7.16, attributed to J. Korevaar, \url{https://arxiv.org/abs/1809.07200}. The record identifies extraction from public arXiv source TeX, accessed July 29, 2026; the 2018 update was retained for status triage.
\item A. Pfluger, cited for the meaning of measurable zero distribution, without a complete bibliographic entry or definition. No authoritative source for the quoted asymptotically analytic growth property is separately supplied.
\item The supplied research summary reports openness as of the 2018 edition and no definitive resolution located, with no additional citation or precise search date.
\end{itemize}
\end{document}
